\subsection{Complexidade O(n)}
Uma complexidade $O(n)$ (complexidade linear) significa que o tempo ou espaço necessário para executar um algoritmo aumenta de forma linear com o tamanho da entrada. Em outras palavras, se você dobrar o tamanho da entrada, o algoritmo também dobrará em tempo de execução ou espaço de memória.

Para demonstrar a complexidade $O(n)$ um exercício retirado do site~\cite{beecrowd1} de busca em um vetor foi escolhido, a atividade em questão consiste em procurar dentro de uma lista de números qual é o menor valor entre eles, destacando também a sua posição dentro da lista.

\begin{center}
\large Menor e Posição

{ BeeCrowd 1180}

Timelimit: $1$	
\end{center}

Faça um programa que leia um valor $N$. Este $N$ será o tamanho de um vetor $X[N]$. A seguir, leia cada um dos valores de X, encontre o menor elemento deste vetor e a sua posição dentro do vetor, mostrando esta informação.

\vspace{0.3cm}
\textbf{Entrada}

A primeira linha de entrada contém um único inteiro $N (1 < N < 1000)$, indicando o número de elementos que deverão ser lidos em seguida para o vetor $X[N]$ de inteiros. A segunda linha contém cada um dos $N$ valores, separados por um espaço. Vale lembrar que nenhuma entrada haverá números repetidos.

\vspace{0.3cm}
\textbf{Saída}
A primeira linha apresenta a mensagem “Menor valor:” seguida de um espaço e do menor valor lido na entrada. A segunda linha apresenta a mensagem “Posição:” seguido de um espaço e da posição do vetor na qual se encontra o menor valor lido, lembrando que o vetor inicia na posição zero.

\begin{table}[!ht]
    \centering
    \begin{tabular}{|p{7cm}|p{7cm}|}
        \hline
        \begin{center}
            \vspace{-0.3cm}
            \textbf{Exemplos de Entrada}
            \vspace{-0.3cm}
        \end{center} 
         &  
        \begin{center}
            \vspace{-0.3cm}
            \textbf{Exemplos de Saída}
            \vspace{-0.3cm}
        \end{center} \\ \hline
         \texttt{10} & \texttt{Menor valor: -5} \\
         \hline
         \texttt{1,2,3,4,-5,6,7,8,9,10} & \texttt{Posicao: 4} \\
         \hline
    \end{tabular}
    \caption{Exemplos de entradas e saídas}
    \label{tabela02}
\end{table}



\begin{table}[ht]
\centering
\caption{Exemplos de entradas e saídas}
\label{tab:algoritmo2}
\begin{tabular}{|p{8cm}|c|c|}
\hline
\textbf{\cellcolor[HTML]{C0C0C0}Algoritmo} & \textbf{\cellcolor[HTML]{C0C0C0}Custo} & \textbf{\cellcolor[HTML]{C0C0C0}Complexidade} \\
\hline
\begin{minipage}{8cm}
\begin{algorithm}[H]
 \caption{Menor e Posicao}
\SetAlgoLined
Declare $X, N, M, POS$ inteiro\;
Leia $N$\;
\For{$X \leftarrow 0$ \textbf{até} $N-1$ faça}{
    Leia $P[X]$\;
}
$M \leftarrow 0$\;
\For{$X \leftarrow 1$ \textbf{até} $N-1$ faça}{
    \If{$P[X] < P[M]$}{
        $M \leftarrow X$\;
    }
}
Escreva $P[M]$\;
Escreva $M$\;
\end{algorithm}
\end{minipage} 
& 
\begin{minipage}{2cm}
C1 \\
C2\\
C3 \\
C4 \\
C5 \\
C6 \\
C7 \\
C8 \\
C9 \\
    C10 \\
C11 \\
C12 \\
C13
\end{minipage}
&
\begin{minipage}{2cm}
1 \\
1 \\
1 \\
1 \\
$n$ \\
$n$ \\
1 \\
$n-1$ \\
$n-1$ \\
$n-1$ \\
$n-1$ \\
1 \\
1
\end{minipage}
\\
\hline
\end{tabular}
\end{table}


\begin{equation}
T(n) = C_1*1 + C_2*1 + C_3*1 + C_4*(n+1) + C_5*n + C_6*1+ C_7*1 + C_8*n
\end{equation}
\begin{equation}
+ C_9*(n-1) + C_10*(n-1) + C_11*1 + C_12*1 + C_13*1 +  C_14*1 + C_15*1
\end{equation}
\begin{equation}
    T(n)= K*(1 + 1 + 1 + n+1 + n + n + 1 + n + n-1 + n-1 + n-1 + n-1 + 1 + 1 + 1)
    \end{equation}
    \begin{equation}
    T(n)= 8*n + 4 =O(n)
    \end{equation}
\begin{center}
\begin{minipage}{0.5\textwidth}
    \lstinputlisting[language=C++]{code/resol++.cpp}
\end{minipage}
\end{center}


    
    \newpage
